\item \emph{A fetch-time survivor panel.} The 470 skips were repositories that no longer resolved publicly in
  September 2026. Equal weighting describes the panel of repositories still public with complete histories,
  not all repositories created in the window.
\item \emph{Surviving stars, not gross additions.} The endpoint reports stars that still exist at retrieval.
  In the 2024 validation the endpoint total was below the archive total for 17 of 20 repositories; for the
  15 large projects the ratio ran from 0.728 to 0.994, with 13 of 15 at least 5 percent below. The study
  target is the surviving-star history; no correction is applied, and later unstars can change the apparent
  past.
\item \emph{Los Angeles day bins.} The endpoint's day boundary was resolved to America/Los\_Angeles days by
  validation, not documented by GitHub; the validation also showed that the transition days around daylight
  saving do not uniformly favor one calendar.
\item \emph{Three windows per repository, and horizons sharing targets.} Inference is conditional on the
  fixed calendar of three cutoffs; the 7- and 30-day targets are prefixes of the 90-day targets, so the
  nine cells per reference are not independent.
\item \emph{Pointwise intervals and BH.} The bootstrap intervals are pointwise, not simultaneous, and the
  Benjamini-Hochberg adjustment carries its usual dependence assumptions.
\item \emph{The signed-rank symmetry assumption.} The Wilcoxon test's location interpretation requires
  symmetric differences, which these are not; it is reported as registered and is not a test of the mean.
\item \emph{AutoETS as a weak uncertainty reference.} Its 80 percent intervals cover about 99 percent of
  target days; a model can beat it on nearly every repository by being less wide. Read with the coverage
  results, the comparison against the local mean is the more informative of the two; that is a reading made
  after the results, and both comparisons stay reported as registered.
\item \emph{Training contamination unresolved.} No author-disclosed complete-corpus training cutoff exists
  for any of the three models; the model cards name specific sources with dates, which do not bound the
  whole corpus. What can be said is bound by availability, per model. The targets, 2025-12-05 to 2026-08-31,
  postdate the pinned TimesFM 2.5 revision, dated 2025-10-02, but are not shown to postdate the pinned
  Chronos-2 revision, dated 2026-06-05, or the pinned TimesFM 3 revision, dated 2026-09-02 with weights
  uploaded 2026-08-24; these are repository revision dates, not training cutoffs. Every repository's history
  before the cutoffs may be familiar to all three models, and a later test window on familiar entities does
  not establish generalization to unseen ones \citep{moghadasi2026}. The controls are Study 1's registered
  follow-up cohort and the lapsed October holdout; neither has run (Section~\ref{sec:notrun}).
\item \emph{The mature-phase calendar.} Every repository is forecast between nine and 41 months after
  creation. Nothing here concerns launches or the days after a release.
\item \emph{Version 1 not recomputed.} The version 1 numbers are quoted from preserved documents.
